\chapter{The 133-Run Core Matrix Study}
\label{chap:133_results}

The 113-Run canonical study plus the 20 descriptive baseline derivations forms the exhaustive backbone of the agricultural decision optimization framework. This section summarizes the statistical derivations separating the best-in-class performing implementations from underperforming strategies.

\section{Overall Leaderboard and Grouped Rankings}

The statistical clustering of the experimental data pinpointed explicit topological ''winners'' within the policy landscape. Figures generated from the canonical \texttt{final\_113} exports heavily indicate that structural choices vastly outplay pure random seed variations.

\begin{figure}[H]
    \centering
    \includegraphics[width=0.9\textwidth]{figures/final_113__leaderboard_primary_metric.png}
    \caption{The Leaderboard showcasing the ranking of primary metrics across the 133 main experiments.}
    \label{fig:113_leaderboard}
\end{figure}

\subsection{Fertilization Results}
In the fertilization framework, the grouped data revealed the \textbf{A2C | nonadaptive | fixed\_weather | years=5000} setting as the paramount grouping. This model secured a mean deterministic return of \textbf{779,267.35} with a 95\% confidence interval of \([727686.53, 830848.18]\). 

Interestingly, at the grouped level, A2C operated most efficiently on long-budget horizons under fixed climate modes. For absolute single-seed peak performance, however, the crown was taken by \textbf{Fertilization | PPO | adaptive | fixed\_weather | years=3000 | seed=1}, hitting a single-run deterministic return of \textbf{791,255.50}.

\begin{figure}[H]
    \centering
    \includegraphics[width=0.9\textwidth]{figures/final_113__grouped_comparison.png}
    \caption{A grouped comparison illustrating variance and mean returns localized by algorithm and weather setup.}
    \label{fig:113_grouped_comp}
\end{figure}

\subsection{Crop-Planning Outcomes}
For standard, non-hierarchical crop planning, \textbf{PPO | nonadaptive | fixed\_weather} triumphed as the best grouped class, recording an unmatched mean evaluation reward of \textbf{21,230.912}.
Conversely, the absolute best single agent iteration occurred with \textbf{A2C | adaptive | fixed\_weather | seed=2}, attaining a peak score of \textbf{23,601.588}. 
The disparity between grouped means and absolute single-run maximums indicates that while PPO presents a much lower variance, high-integrity baseline (making it preferable for deployment), A2C occasionally catches highly optimistic exploitation paths in specific seeds.

\begin{figure}[H]
    \centering
    \includegraphics[width=0.9\textwidth]{figures/final_113__runtime_comparison.png}
    \caption{Runtime Comparison distributions marking the computational efficiency between architectures.}
    \label{fig:113_runtime}
\end{figure}

\section{Hierarchical and Guarded Reruns}

The introduction of \textit{Guarded Hierarchical Reruns} aggressively fundamentally changed the interaction landscape. Due to external constraints, these outcomes demand isolated interpretation. The premier hierarchical group was \textbf{PPO | fixed\_weather | guarded\_rerun}, delivering an astounding mean deterministic return of \textbf{1,861,223.08}.
The single top-performing run across the entire thesis, \textbf{Crop planning hierarchical | PPO | fixed\_weather | seed=1 | profile=pakistan\_baseline}, logged \textbf{1,971,988.25}.

\begin{figure}[H]
    \centering
    \includegraphics[width=0.9\textwidth]{figures/final_113__uplift_vs_baseline.png}
    \caption{Uplift versus Descriptive Baselines: Measuring the net value provided by Deep RL relative to rigid, programmatic agricultural conventions.}
    \label{fig:113_uplift}
\end{figure}

\section{Summary of Core Interpretations}
\begin{enumerate}
    \item The fertilization matrix explicitly favored A2C over the extended 5000-year fixed weather regimes, reinforcing A2C's capacity to maintain stable gradients when environmental noise is artificially bounded.
    \item PPO generated far more consistent behavior across different seeds in typical crop planning, overcoming the brittleness usually associated with generic Actor-Critic gradients.
    \item DQN iterations were empirically relegated to a descriptive baseline tier, routinely outpaced by bounded policy implementations.
\end{enumerate}
